\section*{Author Section 6: derandomizing the multi-switching lemma}
In this section, we prove Lemma 9. The rest of the section is devoted to the proof of Lemma 9. Section 6.1 states some preliminary tools and includes a proof overview. Section 6.2 considers the case that $x$ is truly random. It extends the idea in Section 5.2 and has a rather similar structure. Finally in Section 6.3, we consider the case that $x$ is pseudorandom, and finish the proof.
\subsection*{6.1 Preliminaries}
\paragraph{The canonical partial decision tree.} Let $\rho$ be a random restriction under which $\mathcal F|_\rho$ fails to have a $w$-partial depth-$t$ decision tree. Let $\beta\in\{0,1\}^n$ be an unknown string that our decision tree shall query. We consider the following attempt (see Algorithm 3) to construct a partial decision tree for $\mathcal F|_\rho$. Here $z\in\{0,1\}^n$ is an auxiliary string to be chosen. Every choice of $z\in\{0,1\}^n$ yields a different partial decision tree for $\mathcal F$. In [25], the construction was called ``canonical partial decision tree''.
\paragraph{Algorithm 3: Canonical Partial Decision Tree.} Input: A list of DNFs $\mathcal F=\{F_1,\ldots,F_m\}$, black-box access to a string $\beta\in\{0,1\}^n$, and an auxiliary string $z\in\{0,1\}^n$.
\begin{quote}\textbf{initialize:} $x\leftarrow(\star)^n$; $j\leftarrow1$; $\mathrm{counter}\leftarrow0$.\\
\textbf{while} $\mathrm{counter}<t$ \textbf{do}\\
\hspace*{1em}Find the smallest $i\ge j$ such that $\mathrm{DT}(F_i|_x)>w$. If no such $i$ exists, exit the loop.\\
\hspace*{1em}$y\leftarrow(\star)^n$; $I\leftarrow\emptyset$.\\
\hspace*{1em}\textbf{while} $F_i|_{x\circ y}(\star)$ is not constant and $\mathrm{counter}<t$ \textbf{do}\\
\hspace*{2em}$C_{i,q}\leftarrow$ the term that $T_{F_i|_{x\circ y}}$ will query.\\
\hspace*{2em}$B_{i,q}\leftarrow$ the set of unknown variables in $C_{i,q}|_{x\circ y}$.\\
\hspace*{2em}$y_{B_{i,q}}\leftarrow z_{B_{i,q}}$.\\
\hspace*{2em}$I\leftarrow I\cup B_{i,q}$.\\
\hspace*{2em}$\mathrm{counter}\leftarrow\mathrm{counter}+|B_{i,q}|$.\\
\hspace*{1em}Query $\beta_I$, and set $x_I\leftarrow\beta_I$.\\
\hspace*{1em}$j\leftarrow i$.\\
\textbf{return} $x$.\end{quote}
Note that in Algorithm 3, it is possible that for a formula $F_i$ to get picked by the outer ``while'' loop more than once.
\paragraph{Proof overview.} Before we dive into the formal proof, we try to give some (over-simplified) intuition here. If $\mathcal F$ does not have $w$-partial DT of depth $t$, then Algorithm 3 fails to construct a partial decision tree of depth $t$ on every auxiliary string $z$. However, we will only consider some ``adversarially chosen'' $z$'s. That is, we only consider those $z$'s that trick Algorithm 3 to make at least $w$ queries on each chosen formula. By doing so, we are guaranteed that Algorithm 3 only chooses $t/w$ formulae in total. Then, we can pay a factor of $m^{t/w}$ to enumerate a subset of $t/w$ formulae in $\mathcal F$, and calculate the probability that Algorithm 3 ``gets stuck'' on those formulae. Having fixed those $t/w$ formulae, we can use ideas similar to Section 5.2, and bound the probability by $O(pk)^t$. In summary, we get $\Pr[\text{bad restriction}]\le m^{t/w}O(pk)^t$. Here we only considered the case that the restriction string $x$ is purely random. For the case that both $\Lambda$ and $x$ are pseudorandom, we use tricks similar to Section 5.3 and pay another additive factor.
\subsection*{6.2 Proof when $x$ is truly random}
We start the proof by considering the case that $x$ is a truly random string. Fix $\rho$ to be a bad restriction, under which $\mathcal F|_\rho$ fails to have a $w$-partial depth-$t$ decision tree. It implies that on every $z\in\{0,1\}^n$, Algorithm 3 fails to construct a partial decision tree for $\mathcal F$. However, for ease of our analysis, we will only consider a special class of string $z$. We give the following definition.
\paragraph{Definition 16.} Let $\rho$ be a bad restriction for $\mathcal F$. We call $z\in\{0,1\}^n$ a powerful $(w,t)$-refutation for $\rho$, if there exists $\beta\in\{0,1\}^n$ satisfying the following: Algorithm 3 on input $(\mathcal F|_\rho,\beta,z)$ makes at least $t$ queries. Moreover, at each time we query $\beta_I$, it is guaranteed that $|I|\ge w$.
Powerful refutations always exist for bad restrictions, as shown in the following.
\paragraph{Lemma 17.} If $\mathcal F|_\rho$ does not have $w$-partial decision tree of depth $t$, then there exists a powerful $(w,t)$-refutation for $\rho$.
\begin{proof}If $\mathcal F_\rho$ does not have $w$-partial decision tree of depth $t$, then it means Algorithm 3 fails to construct a $w$-partial DT of depth $t$ for every $z\in\{0,1\}^n$. Then we construct a powerful refutation $z$, together with its associated adversarial input $\beta$, in the following way.
\begin{itemize}\item Initially, we set $z=\beta=(\star)^n$. Then we monitor the execution of Algorithm 3 on $(\mathcal F,\beta,z)$, and gradually fill in $z,\beta$ when they are accessed by Algorithm 3.
\item Whenever Algorithm 3 chooses one formula $F_i|_{\rho\circ x}$ in the outer while-loop, we do the following:
\begin{itemize}\item When running inside the inner while-loop, Algorithm 3 needs to consult the auxiliary string $z$. We can fill in relevant variables of $z$ in such a way that Algorithm 3 consults at least $w$ bits from $z$ in the loop. Since $\mathrm{DT}(F_i|_{\rho\circ x})\ge w$, this is always possible.
\item After finishing the inner loop, Algorithm 3 will make a query to $\beta_I$. At this point, we set $\beta_I$ in such a way that $\mathcal F|_{\rho\circ x\circ\beta}$ does not have $w$-partial DT of depth $(t-\mathrm{counter})$. This is always possible as long as $\mathcal F_{\rho\circ x}$ does not have $w$-partial DT of depth $(t-\mathrm{counter}+|I|)$, which holds by induction.\end{itemize}
\item After Algorithm 3 returns, we fill in the remaining bits of $\beta,z$ with zeros.\end{itemize}\end{proof}
Having established Lemma 17, we come up with the following definition of ``global witness'' naturally.
\paragraph{Definition 18.} Let $t,w$ be two integers. Consider a list of $k$-DNFs $\mathcal F=\{F_1,\ldots,F_m\}$. Suppose $\rho\in\{0,1,\star\}^n$ is a restriction. Let $(R,L_i,S_i,W_i,\beta_i)$ be a tuple, where
\begin{enumerate}\item $1\le R\le t/w$ is an integer;
\item $1\le L_1\le L_2\le\cdots\le L_R\le m$ is a list of $R$ non-decreasing indices;
\item $S_1,\ldots,S_R$ is a list of $R$ integers such that $\sum_{i=1}^RS_i\in[t,t+k]$;
\item $W_1,\ldots,W_R$ is a list of witnesses (as per Definition 12). For every $i\in[R]$, $W_i$ has size $S_i$;
\item $\beta_1,\ldots,\beta_R$ are $R$ strings where $|\beta_i|=S_i$ for every $i\in[R]$.\end{enumerate}
We call the tuple a $(w,t)$-global witness for $\rho$, if it satisfies the following.
\begin{enumerate}\item Set $\rho_1=\rho$. $W_1$ is a $S_1$-witness for $F_{L_1}|_{\rho_1}$.
\item For every $i\ge2$, let $I_{i-1}\subseteq[n]$ be the set of variables involved in $W_{i-1}$. Note that $|I_{i-1}|=S_{i-1}$ since the size of $W_{i-1}$ is $S_{i-1}$. Identify $\beta_{i-1}$ as a partial assignment in $\{0,1,\star\}^n$ where only the part $\beta_{i-1,I_{i-1}}$ is set and other coordinates are filled in with $\star$. Construct $\rho_i=\rho_{i-1}\circ\beta_{i-1}$. Then $W_i$ is a $S_i$-witness for $F_{L_i}|_{\rho_i}$.\end{enumerate}
The size of the global witness is defined as $\sum_{i=1}^RS_i$. As a corollary of Lemma 17, we have:
\paragraph{Corollary 19.} Consider a list of $k$-DNFs $\mathcal F=\{F_1,\ldots,F_m\}$. Suppose $\rho\in\{0,1,\star\}^n$ is a restriction such that $\mathcal F|_\rho$ does not have $w$-partial depth-$t$ decision tree. Then there exists a $(w,t)$-global witness for $\rho$.
\begin{proof}By Lemma 17, there exists a powerful $(w,t)$-refutation for $\mathcal F|_\rho$. Take one such refutation $z$ with its adversarial input $\beta$. Inspect the execution of Algorithm 3 on $(\mathcal F|_\rho,\beta,z)$ and record the transcript. The transcript contains the desired tuple.\end{proof}
Similar to what we have done in Section 5.2, we define the global partial witness.
\paragraph{Definition 20.} Let $t,w$ be two integers. Consider a list of $k$-DNFs $\mathcal F=\{F_1,\ldots,F_m\}$. Suppose $\rho\in\{0,1,\star\}^n$ is a restriction. Let $(R,L_i,S_i,P_i,\beta_i)$ be a tuple, where
\begin{enumerate}\item $1\le R\le t/w$ is an integer;
\item $1\le L_1\le L_2\le\cdots\le L_R\le m$ is a list of $R$ non-decreasing indices;
\item $S_1,\ldots,S_R$ is a list of $R$ integers such that $\sum_{i=1}^RS_i\in[t,t+k]$;
\item $P_1,\ldots,P_R$ is a list of partial witnesses. For every $i\in[R]$, $P_i$ has size $S_i$.
\item $\beta_1,\ldots,\beta_R$ are $R$ strings where $|\beta_i|=S_i$ for every $i\in[R]$.\end{enumerate}
We call $(R,L_i,S_i,P_i,\beta_i)$ a $(w,t)$-global partial witness for $\rho$, if we can complete $P_i$ to get a witness $W_i$ for every $i\in[R]$, such that $(R,L_i,S_i,W_i,\beta_i)$ is a global witness for $\rho$.
\paragraph{Remark 21.} Again, it is easy to see by induction that, for every global partial witness, there is exactly one way to complete it and get a global witness. By a simple counting, it turns out for every $S\in[t,t+k]$, there are at most $2m^{t/w}(12k)^S$ possible global partial witnesses of size $S$:
\begin{enumerate}\item First, $2m^{t/w}$ ways to choose at most $t/w$ formulae in $\mathcal F$.
\item Then, $3^S$ ways to partition the $s$ units of ``query budget'' into $R$ partial witnesses as $S=\sum_{i=1}^RS_i$, and further partition the budget for each partial witness as $S_i=\sum_{j=1}^{r_i}s_{i,j}$.
\item Then, at most $k^S$ ways to construct sets $B_{i,j}$ for each partial witness.
\item Finally, there are $4^S$ ways to choose $\beta_i$'s, as well as $\alpha_{i,j}$'s in each partial witness.\end{enumerate}
Fixing a global partial witness $(R,L_i,S_i,P_i,\beta_i)$, we try to bound the probability that a random $\rho$ has this witness by $(2p)^k$. Using the witness searcher (Algorithm 2) as a subroutine, we design a global witness searcher first. Again, the global witness searcher needs access to an advice string $y\in\{0,1\}^n$. See Algorithm 4. In the following, we use $S$ to denote Algorithm 4 for brevity. Based on Lemma 15, we prove:
\paragraph{Lemma 22.} If $(R,L_i,S_i,P_i,\beta_i)$ is indeed a global partial witness for $\rho$, then there exists an advice $y$ that makes $S$ find $(R,L_i,S_i,P_i,\beta_i)$. More importantly, on a uniformly random $y\sim U_n$, $S$ finds $(\ell_1,\ldots,\ell_r)$ with probability exactly $2^{-S}$, where $S:=\sum_{i=1}^RS_i$ is the size of the global witness.
\paragraph{Algorithm 4: Global Witness Searcher.} Input: A list of DNFs $\mathcal F=\{F_1,\ldots,F_m\}$, a restriction $\rho\in\{0,1,\star\}^n$, a global partial witness $(R,L_i,S_i,P_i,\beta_i)$, and an advice $y\in\{0,1\}^n$.
\begin{quote}\textbf{initialize:} $z\leftarrow\rho\circ y$; $c\leftarrow1$; $\rho^{(1)}\leftarrow\rho$.\\
\textbf{while} $c\le R$ \textbf{do}\\
\hspace*{1em}Run Algorithm 2 on $(F_{L_c},\rho^{(c)},P_c,y)$. If it reports ERROR, report ERROR and terminate the procedure. Otherwise let $W_c$ be the witness returned.\\
\hspace*{1em}$I_c\leftarrow$ the set of variables involved in $W_c$.\\
\hspace*{1em}Identify $\beta_c$ as a partial assignment, where only $\beta_{I_c}$ is fixed.\\
\hspace*{1em}$\rho^{(c+1)}\leftarrow\rho^{(c)}\circ\beta_c$.\\
\hspace*{1em}$c\leftarrow c+1$.\\
\textbf{return} $x$.\end{quote}
\begin{proof}We examine the execution of $S$. It first runs Algorithm 2 as a subroutine, trying to find $W_1$ for $P_1$. With probability $2^{-S_1}$ over $y$, Algorithm 2 succeeds in finding $W_1$. Conditioning on this happened, we have only committed the assignment of $y_{I_1}$, and the rest part of $y$ is still uniformly random. Also note that we set $\rho_{I_1}^{(2)}$ to $\beta_1$, which will hide $y_{I_1}$ in the later execution. This enables us to do an induction and finish the proof.\end{proof}
\paragraph{Decoupling.} Now, we can observe that $S$ only needs to know $z=\rho\circ y$ to work. In particular, it does not need to know which part of $z$ is fixed in $\rho$. Therefore, we can revise $S$ to get a searcher $S'$: the input to $S'$ is now a string $z$ and a global partial witness $(R,L_i,S_i,P_i,\beta_i)$. Otherwise it runs identically the same way as $S$. We denote the output of $S'$ as $S'(z,(R,L_i,S_i,P_i,\beta_i))$. By Lemma 22, we have
\begin{multline*}\Pr_{\rho\sim R}[(R,L_i,S_i,P_i,\beta_i)\text{ is a global partial witness for }\rho]\\\le2^S\Pr_{\rho\sim R,y\sim U_n}[S'(\rho\circ y,(R,L_i,S_i,P_i,\beta_i))\text{ is a global witness for }\rho].\end{multline*}
We observe that
\[\mathbf1_{\{(R,L_i,S_i,W_i,\beta_i)\text{ is a global witness for }\rho\}}\le\mathbf1_{\{(\bigcup_{j=1}^RI_j)\subseteq\Lambda\}}.\tag{11}\]
For a $(t+k)$-wise $p$-bounded $\Lambda$, the event on the right hand side holds with probability $p^{\sum_j|I_j|}$. Then, we have
\begin{align*}&\Pr_{\rho\sim\Lambda[x,0],y\sim U_n}[S'(\rho\circ y,(R,L_i,S_i,P_i,\beta_i))\text{ is a global witness for }\rho]\\&=\mathbb E_\Lambda\mathbb E_{x,y\sim U_n}[\mathbf1_{\{S'(\Lambda[x,y],(R,L_i,S_i,P_i,\beta_i))\text{ is a global witness for }\rho\}}]\\&=\mathbb E_{z\sim U_n}\left[\Pr_\Lambda[S'(z,(R,L_i,S_i,P_i,\beta_i))\text{ is a global witness for }\rho]\right]\\&\le p^S,\end{align*}
where the third line is due to that $\Lambda[x,y]$ is distributed as $U_n$ when $x,y\sim U_n$, and the last inequality holds by (11) and the $(t+k)$-wise $p$-bounded property of $\Lambda$.
\paragraph{Wrapping-up.} We finish the proof by enumerating all $(R,L_i,S_i,P_i,\beta_i)$ and taking a summation.
\begin{align*}&\Pr_{\rho\sim\Lambda[x,\star]}[\mathcal F|_\rho\text{ does not have }w\text{-partial depth-}t\text{ decision tree}]\\&\le\sum_{(R,L_i,S_i,P_i,\beta_i)}\Pr_{\rho\sim\Lambda[x,\star]}[(R,L_i,S_i,P_i,\beta_i)\text{ is a global partial witness for }\rho]\\&\le\sum_{(R,L_i,S_i,P_i,\beta_i)}2^S\mathbb E_{z\in U_n}\left[\Pr_\Lambda[S'(z,(R,L_i,S_i,P_i,\beta_i))\text{ is a global witness for }\Lambda[x,\star]]\right]\\&\le\sum_{S=t}^{t+k}2m^{t/w}2^S(12k)^Sp^S\\&\le4m^{t/w}(24kp)^t.\tag{12}\end{align*}
\subsection*{6.3 Proof when $x$ is pseudorandom}
Now we consider the case that $x$ is not truly random. Our only requirement for $x$ is that, for every fixed $\Lambda$, $x$ is a pseudorandom string that $\epsilon$-fools CNFs of size at most $m$. First of all, by Remark 21, the following equation is established.
\begin{multline*}\mathbf1_{\{(R,L_i,S_i,P_i,\beta_i)\text{ is a global partial witness for }\rho\}}\\=\sum_{W_i:\text{completion of }P_i}\mathbf1_{\{(R,L_i,S_i,W_i,\beta_i)\text{ is a global witness for }\rho\}}.\tag{13}\end{multline*}
Then, the deduction in (12) implies that
\[\sum_{(R,L_i,S_i,W_i,\beta_i)}\Pr_{\rho\sim\Lambda[U_n,\star]}[(R,L_i,S_i,W_i,\beta_i)\text{ is a global witness for }\rho]\le4m^{t/w}(24kp)^t.\tag{14}\]
Next, fixing a tuple $(R,L_i,S_i,W_i,\beta_i)$ and an instantiation of $\Lambda$, we consider the predicate
\[h_\Lambda^{(R,L_i,S_i,W_i,\beta_i)}(x):=\mathbf1_{\{(R,L_i,S_i,W_i,\beta_i)\text{ is a global witness for }\Lambda[x,\star]\}}.\]
We omit the superscript and subscript of $h$ when they are clear from context. We claim that $h$ can be implemented by a CNF of size at most $\sum_{i=1}^m\operatorname{size}(F_i)\le m^2$. To see this, note that $h(x)=1$ if for every $i\in[R]$, $W_i$ is a witness for $\rho\circ\beta_1\circ\cdots\circ\beta_{i-1}$. By argument in Section 6.3, this can be verified using a CNF of size $\operatorname{size}(F_{L_i})$. Therefore, we can evaluate $h$ using an AND over $R$ CNFs, which is itself a larger CNF. Therefore, we conclude that $x$ $\epsilon$-fools $h$. That is,
\[\Pr_{x\sim U_n}[h_\Lambda^{(R,L_i,S_i,W_i,\beta_i)}(x)]-\Pr_{x\sim\boldsymbol x}[h_\Lambda^{(R,L_i,S_i,W_i,\beta_i)}(x)]\le\epsilon.\tag{15}\]
Finally, we have
\begin{align*}&\Pr_{\rho\sim\Lambda[\boldsymbol x,\star]}[\mathcal F|_\rho\text{ does not have }w\text{-partial depth-}t\text{ decision tree}]\\&\le\sum_{(R,L_i,S_i,W_i,\beta_i)}\Pr_{\rho\sim\Lambda[\boldsymbol x,\star]}[(R,L_i,S_i,W_i,\beta_i)\text{ is a global witness for }\rho]\\&\le\sum_{(R,L_i,S_i,W_i,\beta_i)}\epsilon+\Pr_{\rho\sim\Lambda[U_n,\star]}[(R,L_i,S_i,W_i,\beta_i)\text{ is a global witness for }\rho]\\&\le4m^{t/w}(24kp)^t+(24m)^{t+k}\epsilon.\end{align*}
Here, the first line is due to the argument in Section 6.2 and observation (13). The second line is due to (15). The last line holds because (14) and there are at most $(24m)^{t+k}$ possible tuples $(R,L_i,S_i,W_i,\beta_i)$.
